\documentclass[11pt]{article}
\usepackage[margin=1in]{geometry}
\usepackage{amsmath,amssymb}
\usepackage[T1]{fontenc}
\usepackage{lmodern}
\usepackage{hyperref}
\usepackage{array}
\usepackage{longtable}

\title{Notes on the Current RG-Tail Waveform Implementation}
\author{}
\date{\today}

\begin{document}
\maketitle

\section{Purpose}

This note summarizes the current implementation in
\texttt{src/rg\_tail/waveform.py}. It replaces the earlier separate summary,
paper-audit, and \(h(f)\)-prescription notes.

The goal is to describe, in one place:
\begin{itemize}
  \item what the current code implements,
  \item which parts come directly from Section~IV of
  \href{https://arxiv.org/abs/2602.08833}{arXiv:2602.08833},
  \item which parts are repository-level modelling choices,
  \item and what the main current caveats are.
\end{itemize}

\section{High-Level Picture}

The core paper-based object is the factorized \((2,2)\) mode correction
\begin{equation}
  \hat h_{22}
  =
  \hat S_{\rm eff}^{(0)}\,
  T_{22}^{\rm rad}\,
  \rho_{22}^2\,
  e^{i\delta_{22}},
\end{equation}
which is the Section~IV comparable-mass specialization of Eq.~(131) once the
internal-tail bracket is omitted at 4PN.

The current repository then turns this into a practical Fourier-domain
balance-law template by using the dominant-mode flux implied by
\(\hat h_{22}\) inside the adiabatic balance law. The core flux model is
\begin{equation}
  \mathcal F_{22}(x;\lambda_{\rm RG})
  =
  \frac{32}{5}\nu^2 x^5\,\left|\hat h_{22}(x;\lambda_{\rm RG})\right|^2,
\end{equation}
and the Fourier-domain waveform is built from the corresponding
stationary-phase phase evolution.

This distinction is important:
\begin{itemize}
  \item \(\hat h_{22}\) is the paper-facing factorized correction,
  \item \(\mathcal F_{22}\) is the 22-only radiative flux implied by that
  correction,
  \item \(\tilde h(f)\) is the repository's data-analysis wrapper built from
  this flux.
\end{itemize}

\section{What \texttt{waveform.py} Contains}

The module is organized in four blocks:
\begin{enumerate}
  \item test-mass ingredients used for validation,
  \item comparable-mass Section~IV ingredients,
  \item the clean factorized correction \(\hat h_{22}\),
  \item the 22-only balance-law flux and SPA phase helpers,
  \item lightweight detector helpers.
\end{enumerate}

\subsection{Test-mass branch}

The test-mass branch is kept mainly for reproducing the self-convergence
checks discussed in the paper. It includes:
\begin{itemize}
  \item exact Schwarzschild \(E_{\rm circ}\) and \(p_\phi^{\rm circ}\),
  \item the \(\hat \ell_2\) expansion,
  \item test-mass \(\rho_{22}\), \(\delta_{22}\),
  \item test-mass \(\tilde f_{22}\), \(\tilde \delta_{22}\).
\end{itemize}

These are not the main comparable-mass forecast model, but they are useful for
checking that the resummed structure behaves as expected.

\subsection{Comparable-mass Section~IV branch}

For \(\nu > 0\), the module implements:
\begin{itemize}
  \item \(E_{\rm real}/M\) from Eq.~(145),
  \item \(p_{\phi,{\rm circ}}/(\mu M)\) from Eq.~(146),
  \item \(\gamma_{22}^{\rm univ}\) from Eq.~(137),
  \item the radiative tail factor \(T_{22}^{\rm rad}\) from Eq.~(155),
  \item \(\rho_{22}\) from Eq.~(157),
  \item \(\delta_{22}\) from Eq.~(158).
\end{itemize}

The effective source is implemented through
\begin{equation}
  \hat H_{\rm eff}
  =
  \frac{(E_{\rm real}/M)^2 - 1}{2\nu} + 1,
\end{equation}
which is the Eq.~(132) relation solved for the even-parity source.

\section{Explicit Current Model Formulas}

This section collects the main formulas used by the current code in a single
place.

\subsection{Basic variables}

For the dominant \((2,2)\) mode, the code uses
\begin{equation}
  x = (M\Omega)^{2/3},
  \qquad
  \omega = 2\Omega,
  \qquad
  f = \frac{\omega}{2\pi} = \frac{\Omega}{\pi},
  \qquad
  x_f = (\pi M f)^{2/3}.
\end{equation}
The symmetric mass ratio is \(\nu=\eta\).

\subsection{Conservative circular-orbit sector}

The current comparable-mass circular-orbit energy is
\begin{equation}
  \frac{E_{\rm real}}{M}
  =
  1 - \frac{\nu x}{2}
  \left(
    c_0 + c_1 x + c_2 x^2 + c_3 x^3 + c_4 x^4
  \right),
\end{equation}
with
\begin{align}
  c_0 &= 1, \\
  c_1 &= -\frac{3}{4} - \frac{\nu}{12}, \\
  c_2 &= -\frac{27}{8} + \frac{19}{8}\nu - \frac{\nu^2}{24}, \\
  c_3 &=
  -\frac{675}{64}
  + \left(\frac{34445}{576} - \frac{205\pi^2}{96}\right)\nu
  - \frac{155}{96}\nu^2
  - \frac{35}{5184}\nu^3, \\
  c_4 &=
  -\frac{3969}{128}
  + \left(
      -\frac{123671}{5760}
      + \frac{9037\pi^2}{1536}
      + \frac{896\gamma_E}{15}
      + \frac{448}{15}\log(16x)
    \right)\nu
  + \left(
      \frac{498449}{3456}
      - \frac{3157\pi^2}{576}
    \right)\nu^2
  + \frac{301}{1728}\nu^3
  + \frac{77}{31104}\nu^4.
\end{align}

The current circular angular momentum is
\begin{equation}
  \frac{p_{\phi,{\rm circ}}}{\mu M}
  =
  \frac{1}{\sqrt{x}}
  \left(
    d_0 + d_1 x + d_2 x^2 + d_3 x^3 + d_4 x^4
  \right),
\end{equation}
with
\begin{align}
  d_0 &= 1, \\
  d_1 &= \frac{3}{2} + \frac{\nu}{6}, \\
  d_2 &= \frac{27}{8} - \frac{19}{8}\nu + \frac{\nu^2}{24}, \\
  d_3 &=
  \frac{135}{16}
  + \left(
      -\frac{6889}{144} + \frac{41\pi^2}{24}
    \right)\nu
  + \frac{31}{24}\nu^2
  + \frac{7}{1296}\nu^3, \\
  d_4 &=
  \frac{2835}{128}
  + \left(
      \frac{98869}{5760}
      - \frac{128\gamma_E}{3}
      - \frac{6455\pi^2}{1536}
      - \frac{64}{3}\log(16x)
    \right)\nu
  + \left(
      \frac{356035}{3456}
      - \frac{2255\pi^2}{576}
    \right)\nu^2
  - \frac{215}{1728}\nu^3
  - \frac{55}{31104}\nu^4.
\end{align}

The effective source is
\begin{equation}
  \hat H_{\rm eff}
  =
  \frac{(E_{\rm real}/M)^2 - 1}{2\nu} + 1.
\end{equation}

\subsection{RG-tail sector}

The auxiliary quantities entering the universal anomalous dimension are
\begin{equation}
  \hat k = \frac{E_{\rm real}}{M}\,m\Omega,
  \qquad
  J = \frac{p_{\phi,{\rm circ}}/(\mu M)}{(E_{\rm real}/M)^2},
\end{equation}
with \(m=2\) in the current model.

The universal anomalous dimension is
\begin{equation}
  \gamma^{\rm univ}_{22}
  =
  -\frac{214}{105}\hat k^2
  + \frac{2mJ}{3}\hat k^3
  - \frac{3390466}{1157625}\hat k^4
  + \frac{381863\,mJ}{99225}\hat k^5.
\end{equation}

The repository-level RG deformation is
\begin{equation}
  \hat{\hat \ell}_{22}
  =
  2 + \lambda_{\rm RG}\,\gamma^{\rm univ}_{22},
\end{equation}
with GR corresponding to \(\lambda_{\rm RG}=1\).

The radiative-coordinate external tail factor is
\begin{equation}
  T_{22}
  =
  \exp\!\Big[
    \log(120)
    + (\hat{\hat \ell}_{22}-2)\log(2kr_\omega)
    + 2i\hat k \log(4\phi_0)
    + \log\Gamma(\hat{\hat \ell}_{22}-1-2i\hat k)
    - \log\Gamma(2\hat{\hat \ell}_{22}+2)
    + \pi \hat k
    - \frac{i\pi}{2}(\hat{\hat \ell}_{22}-2)
  \Big],
\end{equation}
where
\begin{equation}
  k = 2\Omega,
  \qquad
  r_\omega = \frac{1}{x},
  \qquad
  \phi_0 = \frac{e^{17/12-\gamma_E}}{4}.
\end{equation}

\subsection{Residual amplitude and residual phase}

Define
\begin{equation}
  {\rm eulerlog}_2(x) = \gamma_E + \log(4\sqrt{x}).
\end{equation}
Then the current comparable-mass residual amplitude is
\begin{equation}
  \rho_{22}(x,\nu)
  =
  1 + r_1 x + r_2 x^2 + r_3 x^3 + r_4 x^4,
\end{equation}
with
\begin{align}
  r_1 &= -\frac{43}{42} + \frac{55}{84}\nu, \\
  r_2 &= -\frac{20555}{10584} - \frac{33025}{21168}\nu + \frac{19583}{42336}\nu^2, \\
  r_3 &=
  -\frac{4296031}{4889808}
  + \left(\frac{41\pi^2}{192} - \frac{48993925}{9779616}\right)\nu
  - \frac{6292061}{3259872}\nu^2
  + \frac{10620745}{39118464}\nu^3, \\
  r_4 &=
  \frac{9228174993589}{800950550400}
  + \left(
      -\frac{2487107795131}{145627372800}
      + \frac{464}{35}{\rm eulerlog}_2(x)
      - \frac{9953\pi^2}{21504}
    \right)\nu \nonumber \\
  &\quad
  + \left(
      \frac{10815863492353}{640760440320}
      - \frac{3485\pi^2}{5376}
    \right)\nu^2
  - \frac{2088847783}{11650189824}\nu^3
  + \frac{70134663541}{512608352256}\nu^4.
\end{align}

The current residual phase is
\begin{equation}
  \delta_{22}
  =
  -\frac{17}{3}y^{3/2}
  - 24\nu\,y^{5/2}
  + \left(
      \frac{30995}{1134}\nu + \frac{962}{135}\nu^2
    \right)y^{7/2}
  - \frac{4976\pi}{105}\nu\,y^4,
\end{equation}
where
\begin{equation}
  y = \left( \frac{E_{\rm real}}{M} x^{3/2} \right)^{2/3}.
\end{equation}

The final Section~IV-inspired factorized correction is therefore
\begin{equation}
  \hat h_{22}
  =
  \hat H_{\rm eff}\,T_{22}\,\rho_{22}^2\,e^{i\delta_{22}}.
\end{equation}

\subsection{Flux, balance law, and SPA phase}

Using the Newtonian mode
\begin{equation}
  h_{22}^N = -8\sqrt{\frac{\pi}{5}}\,\nu x\,e^{-2i\phi},
\end{equation}
the current 22-only flux is
\begin{equation}
  \mathcal F_{22}(x;\lambda_{\rm RG})
  =
  \frac{32}{5}\nu^2 x^5
  \left|
    \hat h_{22}(x;\lambda_{\rm RG})
  \right|^2.
\end{equation}

The code uses the balance law
\begin{equation}
  \frac{dt}{dx}
  =
  -M\,
  \frac{d(E_{\rm real}/M)/dx}{\mathcal F_{22}(x)},
\end{equation}
with
\begin{equation}
  \frac{d\Phi_{\rm GW}}{dx}
  =
  2\Omega \frac{dt}{dx}
  =
  2\frac{x^{3/2}}{M}\frac{dt}{dx}.
\end{equation}

If \(x_{\rm ref}\) is the reference endpoint used to anchor the integrals, the
code constructs
\begin{equation}
  \Delta t(x)
  = \int_x^{x_{\rm ref}} \frac{dt}{dx'}\,dx',
  \qquad
  \Delta \Phi_{\rm GW}(x)
  = \int_x^{x_{\rm ref}} \frac{d\Phi_{\rm GW}}{dx'}\,dx',
\end{equation}
and then uses the SPA phase
\begin{equation}
  \Psi_{22,{\rm bal}}(f;\lambda_{\rm RG})
  =
  \Delta \Phi_{\rm GW}(x_f)
  - 2\pi f\,\Delta t(x_f).
\end{equation}

The current default Fourier-domain waveform uses the same balance law for the
SPA amplitude. The Newtonian restricted amplitude is
\begin{equation}
  A_N(f)
  =
  \sqrt{\frac{5}{24}}\,
  \frac{\mathcal M^{5/6}}{d_L\,\pi^{2/3}}\,
  f^{-7/6}.
\end{equation}
Because the time-domain mode amplitude contains
\(|\hat h_{22}|\), while the inspiral rate contains
\(\mathcal F_{22}\propto |\hat h_{22}|^2\), the explicit
\(|\hat h_{22}|\) cancels in the stationary-phase amplitude:
\begin{equation}
  |\hat h_{22}|\,
  \sqrt{\frac{\dot f_N}{\dot f}}
  =
  \sqrt{
    \frac{d(E_{\rm real}/M)/dx}{d(E_N/M)/dx}
  }
  =
  \sqrt{
    -\frac{2}{\nu}\frac{d(E_{\rm real}/M)}{dx}
  }.
\end{equation}
Thus the default detector-domain template is
\begin{equation}
  \tilde h(f)
  =
  A_N(f)\,
  \sqrt{
    -\frac{2}{\nu}\frac{d(E_{\rm real}/M)}{dx}\bigg|_{x_f}
  }\,
  \exp\!\left[
    i\left(
      2\pi f t_c - \phi_c - \frac{\pi}{4}
      + \Psi_{22,{\rm bal}}(f;\lambda_{\rm RG})
      + \arg \hat h_{22}(x_f)
    \right)
  \right],
\end{equation}
with
\begin{equation}
  x_f = (\pi M f)^{2/3}.
\end{equation}
The older project proxy,
\(\tilde h_{\rm proxy}=A_N|\hat h_{22}|\exp(i\Psi)\), is retained in the code
as \texttt{amplitude\_model="proxy"} for controlled comparisons.

\subsection{Code-level \texttt{h\_of\_f} prescription}

The current public waveform entry point is
\texttt{h\_of\_f(f, Mc, eta, dL, tc, phic, lambda\_RG, ...)}. In terms of
the physical masses, the code first converts
\begin{equation}
  M = \mathcal M_c\,\nu^{-3/5},
  \qquad
  d_L = d_L^{\rm Mpc}\,{\rm Mpc},
  \qquad
  M = M^{M_\odot}\,M_\odot,
\end{equation}
with \(G=c=1\). It then defines the Schwarzschild ISCO frequency
\begin{equation}
  f_{\rm ISCO}
  =
  \frac{1}{6^{3/2}\pi M},
  \qquad
  f_{\rm cut}
  =
  \alpha_{\rm cut}\,f_{\rm ISCO},
\end{equation}
where \(\alpha_{\rm cut}=\texttt{fmax\_over\_fisco}\). The waveform is set to
zero outside the code band,
\begin{equation}
  \tilde h(f)=0,
  \qquad
  f<f_{\rm low}
  \quad\hbox{or}\quad
  f\ge f_{\rm cut}.
\end{equation}

Inside the band, the exact current default is
\begin{equation}
  \boxed{
  \tilde h(f)
  =
  W(f)\,
  A_{\rm bal}(f)\,
  \exp\!\left[i\Psi_{\rm code}(f)\right]
  },
\end{equation}
where the Fermi taper is
\begin{equation}
  W(f)
  =
  \frac{1}{
    1+\exp\!\left[
      \frac{f-f_{\rm ISCO}}
      {\sigma_{\rm taper} f_{\rm ISCO}}
    \right]
  },
  \qquad
  \sigma_{\rm taper}
  =
  \texttt{sigma\_taper\_over\_fisco}.
\end{equation}
The code phase is
\begin{equation}
  \Psi_{\rm code}(f)
  =
  2\pi f t_c
  - \phi_c
  - \frac{\pi}{4}
  + \Psi_{22,{\rm bal}}(f;\lambda_{\rm RG})
  + \arg\hat h_{22}(x_f;\lambda_{\rm RG}),
\end{equation}
with \(x_f=(\pi M f)^{2/3}\). The default amplitude,
\texttt{amplitude\_model="balance"}, is
\begin{equation}
  A_{\rm bal}(f)
  =
  A_N(f)
  \sqrt{
    \max\!\left[
      -\frac{2}{\nu}\frac{d(E_{\rm real}/M)}{dx}\bigg|_{x_f},
      0
    \right]
  }.
\end{equation}
The \(\max[\cdots,0]\) is a numerical guard in the code. For the current
fiducial binaries and paper-facing frequency band, the physical branch has
\(d(E_{\rm real}/M)/dx<0\).

The optional amplitude choices are
\begin{align}
  A_{\rm proxy}(f;\lambda_{\rm RG})
  &=
  A_N(f)\,
  \left|\hat h_{22}(x_f;\lambda_{\rm RG})\right|,
  \\
  A_{\rm proxy,phase\ only}(f)
  &=
  A_N(f)\,
  \left|\hat h_{22}(x_f;\lambda_{\rm RG}=1)\right|,
  \\
  A_{\rm newtonian}(f)
  &=
  A_N(f).
\end{align}
Thus \texttt{phase\_only=True} only changes the old proxy amplitude; it does
not change the current default balance-law amplitude, whose explicit
\(\lambda_{\rm RG}\)-dependent magnitude cancels in the SPA construction.

\section{Current Relation to arXiv:2602.08833}

\renewcommand{\arraystretch}{1.15}
\begin{longtable}{|>{\raggedright\arraybackslash}p{0.18\textwidth}|>{\raggedright\arraybackslash}p{0.18\textwidth}|>{\raggedright\arraybackslash}p{0.12\textwidth}|>{\raggedright\arraybackslash}p{0.38\textwidth}|}
\hline
\textbf{Paper item} & \textbf{Code function} & \textbf{Status} & \textbf{Current note} \\
\hline
\endfirsthead
\hline
\textbf{Paper item} & \textbf{Code function} & \textbf{Status} & \textbf{Current note} \\
\hline
\endhead

Eqs.~(85), (86) & \texttt{E\_circ\_test\_mass}, \texttt{p\_phi\_circ\_test\_mass} & Match & Implemented in exact Schwarzschild form. \\
\hline

Eq.~(175) & \texttt{ellhat\_test\_mass} & Match & The explicit test-mass \(\hat \ell_2\) coefficients are kept. \\
\hline

Eq.~(112) & \texttt{rho\_22\_test\_mass} & Match & The current code now keeps the test-mass series through \(x^{10}\). \\
\hline

Eq.~(113) & \texttt{ftilde\_22\_test\_mass} & Match & The current code now keeps the test-mass series through \(x^{10}\). \\
\hline

Eqs.~(114), (115) & \texttt{delta\_22\_test\_mass}, \texttt{deltatilde\_22\_test\_mass} & Match & Implemented as written in the paper. \\
\hline

Eq.~(145) & \texttt{energy\_real\_over\_M} & Match & Implemented at 4PN. \\
\hline

Eq.~(146) & \texttt{p\_phi\_circ\_over\_muM} & Match & Implemented at 4PN. \\
\hline

Eq.~(137) & \texttt{gamma\_univ\_22} & Match & Implemented with the paper's \(\hat k\) and \(J\) definitions. \\
\hline

Eq.~(155) & \texttt{tail\_factor\_T22} & Match up to intended 4PN use & The external radiative tail is implemented in log-space for numerical stability. \\
\hline

Eq.~(157) & \texttt{rho\_22} & Match & Implemented at 4PN. \\
\hline

Eq.~(158) & \texttt{delta\_22} & Match & Implemented at 4PN. \\
\hline

Eq.~(131), 4PN Section~IV use & \texttt{hhat\_22} & Match to intended 4PN model & The external-tail-only 4PN form is used, consistent with the Section~IV discussion below Eq.~(143). \\
\hline

Internal-tail bracket in Eq.~(131) & \texttt{include\_internal\_tail} option & Not implemented & The code now raises \texttt{NotImplementedError} instead of silently pretending to include it. \\
\hline

22-only flux \(\mathcal F_{22}\) & \texttt{flux\_22} & Repository choice consistent with standard multipolar flux construction & Built from \(h_{22}=h_{22}^N \hat h_{22}\) as \(\mathcal F_{22}=(32/5)\nu^2 x^5 |\hat h_{22}|^2\). \\
\hline

Balance-law phase & \texttt{\_spa\_phase\_from\_22\_flux} & Repository choice & The paper does not give a final SPA detector template, so the phase is built by integrating the 22-only flux with the GR conservative energy. \\
\hline

Detector-template \(\tilde h(f)\) & \texttt{h\_of\_f} & Repository choice & The final Fourier-domain waveform uses a balance-law SPA phase and, by default, the SPA amplitude correction implied by the same balance law. The older \(A_N|\hat h_{22}|\) proxy is still available as an option. \\
\hline

\end{longtable}

\section{Comparison with the 4PN Standard \(H_{22}\) Mode}

The current model is also useful to compare against the standard 4PN
dominant-mode result in
\href{https://arxiv.org/abs/2304.11185}{arXiv:2304.11185}. In that paper,
Eq.~(11) gives the complex \((2,2)\) mode \(H_{22}\) through relative 4PN
order, using the normalization
\begin{equation}
  h_+ - i h_\times
  =
  \frac{8Gm\nu x}{Rc^2}\sqrt{\frac{\pi}{5}}
  \sum_{\ell m} H_{\ell m}e^{-im\psi}\,{}_{-2}Y_{\ell m}.
\end{equation}
With this normalization, the corresponding 22-mode flux is
\begin{equation}
  \mathcal F_{22}^{2304.11185}
  =
  \frac{32}{5}\nu^2x^5 |H_{22}|^2.
\end{equation}

For \(\lambda_{\rm RG}=1\), the current code agrees with
Eq.~(11) through relative 3.5PN order. The first difference appears at
relative 4PN:
\begin{equation}
  H_{22}^{2304.11185} - H_{22}^{\rm code}
  =
  \frac{16}{45}(\nu-1)
  \left[
    30\,{\rm eulerlog}_2(x)-137-15i\pi
  \right]x^4,
\end{equation}
where
\begin{equation}
  {\rm eulerlog}_2(x)
  =
  \gamma_E+\ln(4\sqrt{x}).
\end{equation}
Equivalently, at the level of the 22-mode flux,
\begin{equation}
  \frac{
    \mathcal F_{22}^{2304.11185}
    -
    \mathcal F_{22}^{\rm code}
  }{
    (32/5)\nu^2x^5
  }
  =
  \frac{32}{45}(\nu-1)
  \left[
    30\,{\rm eulerlog}_2(x)-137
  \right]x^4.
\end{equation}

This term has the logarithmic plus \(i\pi\) structure characteristic of
hereditary radiative propagation effects. The companion derivation paper,
\href{https://arxiv.org/abs/2304.11186}{arXiv:2304.11186}, identifies a new
4PN ``tail-of-memory'' contribution in the radiative quadrupole moment, along
with related tail-looking interactions. The current code should therefore be
read as the Section~IV external-tail factorized model of arXiv:2602.08833, not
as the exact standard 4PN \(H_{22}\) mode of arXiv:2304.11185.

For now we do not add this 4PN difference to the waveform. If exact agreement
with the standard 4PN \(H_{22}\) mode is desired later, one can multiply the
current \(\hat h_{22}\) by
\begin{equation}
  1+
  \frac{16}{45}(\nu-1)
  \left[
    30\,{\rm eulerlog}_2(x)-137-15i\pi
  \right]x^4,
\end{equation}
or equivalently add the corresponding real flux correction above.

\section{Current Beyond-GR Convention}

The parameter \(\lambda_{\rm RG}\) is not part of the GR waveform in the paper.
It is a repository-level deformation used for Fisher forecasts. The current
convention is
\begin{equation}
  \hat{\hat \ell}_{22}
  =
  2 + \lambda_{\rm RG}\,\gamma^{\rm univ}_{22},
\end{equation}
with GR corresponding to \(\lambda_{\rm RG}=1\).

The same scaling is applied in the test-mass branch by rescaling the running
part \(\hat \ell_2 - 2\).

\section{Current Flux and Fourier-Domain Prescription}

The paper gives the radiative mode structure and the harmonic relation
\(\omega = m \Omega\), but it does not give a final detector-template
stationary-phase waveform. The current code therefore first constructs the
dominant-mode flux from the standard multipolar flux formula,
\begin{equation}
  \mathcal F
  =
  \frac{1}{8\pi}\sum_{\ell=2}^{\infty}\sum_{m=1}^{\ell}
  (m\Omega)^2 |h_{\ell m}|^2,
\end{equation}
truncated to the \((2,2)\) mode only:
\begin{equation}
  \mathcal F_{22}
  =
  \frac{1}{8\pi}(2\Omega)^2 |h_{22}|^2.
\end{equation}

With the Newtonian factorization
\begin{equation}
  h_{22}=h_{22}^N \hat h_{22},
  \qquad
  h_{22}^{N}=-8\sqrt{\frac{\pi}{5}}\,\nu x\,e^{-2i\phi},
\end{equation}
this becomes
\begin{equation}
  \mathcal F_{22}(x;\lambda_{\rm RG})
  =
  \frac{32}{5}\nu^2 x^5
  \left|
    \hat h_{22}(x;\lambda_{\rm RG})
  \right|^2.
\end{equation}

Since
\begin{equation}
  \hat h_{22}
  =
  \hat H_{\rm eff}\,
  T_{22}\,
  \rho_{22}^2\,
  e^{i\delta_{22}},
\end{equation}
the current code is therefore using
\begin{equation}
  \mathcal F_{22}(x;\lambda_{\rm RG})
  =
  \frac{32}{5}\nu^2 x^5\,
  \hat H_{\rm eff}(x,\nu)^2\,
  |T_{22}(x,\nu;\lambda_{\rm RG})|^2\,
  \rho_{22}(x,\nu)^4.
\end{equation}

The SPA phase is then built from the adiabatic balance law
\begin{equation}
  \frac{dt}{dx}
  =
  -M\,
  \frac{d(E_{\rm real}/M)/dx}{\mathcal F_{22}(x)},
\end{equation}
with the circular-orbit relation \(\Omega=x^{3/2}/M\). The Fourier-domain
phase is assembled from the corresponding time-to-reference and
gravitational-wave phase-to-reference integrals. The resulting default signal
model for the \((2,2)\) mode is
\begin{equation}
  \tilde h(f)
  =
  A_N(f)\,
  \sqrt{
    -\frac{2}{\nu}\frac{d(E_{\rm real}/M)}{dx}\bigg|_{x_f}
  }\,
  \exp\!\left[
    i\left(
      2\pi f t_c - \phi_c - \frac{\pi}{4}
      + \Psi_{22,{\rm bal}}(f;\lambda_{\rm RG})
      + \arg \hat h_{22}(x_f)
    \right)
  \right],
\end{equation}
with
\begin{equation}
  x_f = (\pi M f)^{2/3}.
\end{equation}
Here
\begin{equation}
  A_N(f)
  =
  \sqrt{\frac{5}{24}}\,
  \frac{\mathcal M^{5/6}}{d_L\,\pi^{2/3}}\,
  f^{-7/6}.
\end{equation}

So the current default \(h(f)\) is:
\begin{quote}
restricted Newtonian SPA amplitude + conservative-energy SPA amplitude
correction + 22-only balance-law SPA phase + complex Section~IV phase
\(\arg\hat h_{22}(x_f)\).
\end{quote}

This should be viewed as:
\begin{itemize}
  \item paper-based in the construction of \(\hat h_{22}\),
  \item more self-consistent than the older TaylorF2 wrapper because the phase
  is now driven by the same 22-mode flux that carries the RG deformation,
  \item still approximate because it uses only the \((2,2)\) mode and a
  simplified detector-domain projection.
\end{itemize}

\subsection{Amplitude options}

The code keeps the older project proxy for diagnostics:
\begin{equation}
  A_{\rm proxy}(f;\lambda_{\rm RG})
  =
  A_N(f)\,|\hat h_{22}(x_f;\lambda_{\rm RG})|.
\end{equation}
With \texttt{amplitude\_model="proxy"} and \texttt{phase\_only=True}, this
proxy amplitude is frozen to its GR value,
\begin{equation}
  A_{\rm proxy,phase\ only}(f)
  =
  A_N(f)\,|\hat h_{22}(x_f;\lambda_{\rm RG}=1)|,
\end{equation}
while the phase still uses the requested \(\lambda_{\rm RG}\)-dependent
22-only balance-law dynamics and the requested \(\arg \hat h_{22}(x_f)\).
For the default \texttt{amplitude\_model="balance"}, the amplitude has no
explicit \(\lambda_{\rm RG}\)-dependent magnitude after the SPA cancellation,
so \texttt{phase\_only} does not change the default amplitude.

\section{Practical Status}

At the current point, the implementation is best described as:
\begin{itemize}
  \item a clean Section~IV-inspired \(\hat h_{22}\) implementation,
  \item plus a practical 22-only balance-law SPA wrapper for Fisher work,
  \item with the test-mass branch retained for local validation.
\end{itemize}

It is suitable for the current standalone RG-tail Fisher studies, but it is
not yet a full EOB waveform model and it is not integrated into the main
population/network forecast pipeline in the repository.

\section{Current Paper-Facing Figures}

For the systems paper, the RG-tail material should stay compact. The current
recommended figure set is:
\begin{itemize}
  \item an expanded detector-and-network comparison plot of
  \(\sigma(\lambda_{\rm RG})\), generated from
  \texttt{results/rg\_tail\_forecast/summary/fisher\_results.json};
  \item a Gaussian Fisher corner plot for the CE 40 km forecast, showing the
  covariance among \(\mathcal M_c\), \(\eta\), \(d_L\), and
  \(\lambda_{\rm RG}\).
\end{itemize}

The corresponding current files are:
\begin{itemize}
  \item \texttt{results/rg\_tail\_forecast/summary/sigma\_lambda\_rg.pdf},
  \item \texttt{results/rg\_tail\_forecast/summary/rg\_tail\_corner\_ce40.pdf}.
\end{itemize}

The detector-and-network comparison figure is the cleanest headline output for
the skill. In its current paper-facing version it includes:
\begin{equation}
  \text{aLIGO},\quad
  \text{Voyager},\quad
  \text{CE 20 km},\quad
  \text{CE 40 km},\quad
  \text{ET},\quad
  \text{CE 40+20 km},\quad
  \text{ET + CE 40 km},\quad
  \text{ET + CE 40+20 km}.
\end{equation}
For the current GW150914-like fiducial source with
\(f_{\max}=f_{\rm ISCO}\), the current comparison gives
\begin{equation}
  \sigma(\lambda_{\rm RG})
  \approx
  \begin{cases}
    4.28\times 10^{-1}, & \text{aLIGO}, \\
    2.14\times 10^{-2}, & \text{Voyager}, \\
    3.58\times 10^{-3}, & \text{CE 20 km}, \\
    5.29\times 10^{-4}, & \text{CE 40 km}, \\
    6.00\times 10^{-3}, & \text{ET}, \\
    5.02\times 10^{-4}, & \text{CE 40+20 km}, \\
    5.25\times 10^{-4}, & \text{ET + CE 40 km}, \\
    4.98\times 10^{-4}, & \text{ET + CE 40+20 km},
  \end{cases}
\end{equation}
where the network points are reduced-network forecasts obtained by summing the
independent Fisher matrices of the constituent detectors.
In this reduced model, CE 40 km already carries most of the constraining power,
so the network improvements are real but modest.

The CE corner plot is useful because it shows that the present
\(\lambda_{\rm RG}\) constraint is not only a one-number summary; it also
reveals the dominant parameter degeneracies of the current skill-level model,
especially the strong anticorrelation with \(\mathcal M_c\) and the milder
correlations with \(\eta\) and \(d_L\).

\section{Main Caveats}

\begin{itemize}
  \item The internal-tail factor in Eq.~(131) is still omitted.
  \item The \(\lambda_{\rm NS}^{\rm inst}\) denominator is not part of the
  current default tail implementation.
  \item The final \(\tilde h(f)\) is still a repository-level stationary-phase
  construction; arXiv:2602.08833 provides the mode ingredients, not a complete
  detector-template prescription.
  \item The flux keeps only the \((2,2)\) mode.
  \item The conservative dynamics are still the GR circular-orbit energy from
  Eq.~(145); the RG deformation currently enters the radiative sector only,
  through the tail factor in \(\hat h_{22}\).
  \item The network comparison is a reduced Fisher sum of constituent detector
  information and does not yet include sky location, polarization, or
  antenna-pattern geometry.
  \item The SPA amplitude is now derived from the same 22-only balance law, but
  the waveform still omits higher modes, spins, merger-ringdown, and detector
  response geometry.
  \item The whole RG-tail module remains a standalone research workflow rather
  than a production waveform package.
\end{itemize}

\section{Bottom Line}

The most concise current summary is:
\begin{quote}
\texttt{waveform.py} now implements a clean Section~IV external-tail
\(\hat h_{22}\), derives the corresponding 22-only flux
\(\mathcal F_{22}=(32/5)\nu^2 x^5 |\hat h_{22}|^2\), and uses that flux to
build both the current Fourier-domain SPA phase and the default balance-law
SPA amplitude for Fisher forecasts.
\end{quote}

That is the current up-to-date description of the code.

\end{document}
